% Options for packages loaded elsewhere
\PassOptionsToPackage{unicode}{hyperref}
\PassOptionsToPackage{hyphens}{url}
\documentclass[
  10pt,
  a4paper,
]{article}
\usepackage{xcolor}
\usepackage[margin=22mm]{geometry}
\usepackage{amsmath,amssymb}
\setcounter{secnumdepth}{-\maxdimen} % remove section numbering
\usepackage{iftex}
\ifPDFTeX
  \usepackage[T1]{fontenc}
  \usepackage[utf8]{inputenc}
  \usepackage{textcomp} % provide euro and other symbols
\else % if luatex or xetex
  \usepackage{unicode-math} % this also loads fontspec
  \defaultfontfeatures{Scale=MatchLowercase}
  \defaultfontfeatures[\rmfamily]{Ligatures=TeX,Scale=1}
\fi
\usepackage{lmodern}
\ifPDFTeX\else
  % xetex/luatex font selection
\fi
% Use upquote if available, for straight quotes in verbatim environments
\IfFileExists{upquote.sty}{\usepackage{upquote}}{}
\IfFileExists{microtype.sty}{% use microtype if available
  \usepackage[]{microtype}
  \UseMicrotypeSet[protrusion]{basicmath} % disable protrusion for tt fonts
}{}
\makeatletter
\@ifundefined{KOMAClassName}{% if non-KOMA class
  \IfFileExists{parskip.sty}{%
    \usepackage{parskip}
  }{% else
    \setlength{\parindent}{0pt}
    \setlength{\parskip}{6pt plus 2pt minus 1pt}}
}{% if KOMA class
  \KOMAoptions{parskip=half}}
\makeatother
\usepackage{longtable,booktabs,array}
\newcounter{none} % for unnumbered tables
\usepackage{calc} % for calculating minipage widths
% Correct order of tables after \paragraph or \subparagraph
\usepackage{etoolbox}
\makeatletter
\patchcmd\longtable{\par}{\if@noskipsec\mbox{}\fi\par}{}{}
\makeatother
% Allow footnotes in longtable head/foot
\IfFileExists{footnotehyper.sty}{\usepackage{footnotehyper}}{\usepackage{footnote}}
\makesavenoteenv{longtable}
\setlength{\emergencystretch}{3em} % prevent overfull lines
\providecommand{\tightlist}{%
  \setlength{\itemsep}{0pt}\setlength{\parskip}{0pt}}
\usepackage{mathrsfs}
\newcommand{\Log}{\operatorname{Log}}
\usepackage{mathtools}
\usepackage{xurl}
\emergencystretch=3em
\setcounter{tocdepth}{1}
\newcommand{\Beta}{\mathrm{B}}
\DeclareUnicodeCharacter{00B2}{\ensuremath{^2}}
\DeclareUnicodeCharacter{00B1}{\ensuremath{\pm}}
\DeclareUnicodeCharacter{00D7}{\ensuremath{\times}}
\DeclareUnicodeCharacter{2192}{\ensuremath{\to}}
\DeclareUnicodeCharacter{21D2}{\ensuremath{\Rightarrow}}
\DeclareUnicodeCharacter{21A6}{\ensuremath{\mapsto}}
\DeclareUnicodeCharacter{2208}{\ensuremath{\in}}
\DeclareUnicodeCharacter{2209}{\ensuremath{\notin}}
\DeclareUnicodeCharacter{2212}{\ensuremath{-}}
\DeclareUnicodeCharacter{2227}{\ensuremath{\wedge}}
\DeclareUnicodeCharacter{2228}{\ensuremath{\vee}}
\DeclareUnicodeCharacter{2260}{\ensuremath{\ne}}
\DeclareUnicodeCharacter{2264}{\ensuremath{\le}}
\DeclareUnicodeCharacter{2265}{\ensuremath{\ge}}
\DeclareUnicodeCharacter{2297}{\ensuremath{\otimes}}
\DeclareUnicodeCharacter{00B3}{\ensuremath{^3}}
\DeclareUnicodeCharacter{2074}{\ensuremath{^4}}
\DeclareUnicodeCharacter{2076}{\ensuremath{^6}}
\DeclareUnicodeCharacter{03BB}{\ensuremath{\lambda}}
\DeclareUnicodeCharacter{03BC}{\ensuremath{\mu}}
\DeclareUnicodeCharacter{03B5}{\ensuremath{\epsilon}}
\DeclareUnicodeCharacter{03C4}{\ensuremath{\tau}}
\DeclareUnicodeCharacter{03A6}{\ensuremath{\Phi}}
\DeclareUnicodeCharacter{03B1}{\ensuremath{\alpha}}
\DeclareUnicodeCharacter{03C3}{\ensuremath{\sigma}}
\DeclareUnicodeCharacter{03B4}{\ensuremath{\delta}}
\DeclareUnicodeCharacter{03C0}{\ensuremath{\pi}}
\DeclareUnicodeCharacter{039B}{\ensuremath{\Lambda}}
\DeclareUnicodeCharacter{2202}{\ensuremath{\partial}}
\DeclareUnicodeCharacter{0393}{\ensuremath{\Gamma}}
\DeclareUnicodeCharacter{2205}{\ensuremath{\varnothing}}
\DeclareUnicodeCharacter{221E}{\ensuremath{\infty}}
\DeclareUnicodeCharacter{2200}{\ensuremath{\forall}}
\DeclareUnicodeCharacter{00B9}{\ensuremath{^1}}
\DeclareUnicodeCharacter{00B0}{\ensuremath{{}^\circ}}
\DeclareUnicodeCharacter{211D}{\ensuremath{\mathbb{R}}}
\usepackage{bookmark}
\IfFileExists{xurl.sty}{\usepackage{xurl}}{} % add URL line breaks if available
\urlstyle{same}
\hypersetup{
  hidelinks,
  pdfcreator={LaTeX via pandoc}}

\author{}
\date{}

\begin{document}

{
\setcounter{tocdepth}{3}
\tableofcontents
}
\section{Exact coefficient bound for the original prime-power
window}\label{exact-coefficient-bound-for-the-original-prime-power-window}

For the original von Mangoldt function, including every prime power,
define \[
S(a)=\sum_{\substack{n\ge2\\|\log n-a|\le1/32}}
\frac{\Lambda(n)}{\sqrt n},
\qquad \frac{583}{800}\le a\le\log256.
\tag{FPC1}
\] All logarithms are natural. The exact maximum on this entire real
interval is \[
\boxed{
\max_a S(a)=M_*:=
\frac{\log227}{\sqrt{227}}+
\frac{\log229}{\sqrt{229}}+
\frac{\log233}{\sqrt{233}}+
\frac{\log239}{\sqrt{239}}+
\frac{\log241}{\sqrt{241}}.
}
\tag{FPC2}
\] Its certified rational enclosure is \[
\frac{1783796246}{10^9}<M_*<\frac{1783796247}{10^9}
<\boxed{\frac{223}{125}}<2.
\tag{FPC3}
\] This calculation is finite arithmetic. It does not evaluate or
truncate any zeta-zero sum.

\subsection{The exact finite range and ratio
reduction}\label{the-exact-finite-range-and-ratio-reduction}

For \(0<x<1\), comparison of the exponential series with the geometric
series gives \(e^x<1/(1-x)\). Therefore \[
n\le e^{a+1/32}\le256e^{1/32}
<256\frac{32}{31}=\frac{8192}{31}<265.
\tag{FPC4}
\] Thus all contributing indices satisfy \(n\le264\). Both primes 257
and 263 must be retained. In fact \(e^{1/32}>1+1/32\) gives
\(256e^{1/32}>264\), so the exact integer upper cutoff is 264.

The first index 2 never occurs. The elementary logarithm series gives \[
\log2=2\sum_{j\ge0}\frac{(1/3)^{2j+1}}{2j+1}
<\frac23+\frac{2(1/3)^3}{3(1-1/9)}
=\frac{25}{36}<\frac{279}{400}
=\frac{583}{800}-\frac1{32}.
\tag{FPC5}
\] The strict series bound follows because each denominator after the
first term is at least 3, and those after that are strictly larger.
Non-prime-power indices have coefficient zero and are preserved with
that numerical value.

Any two indices \(m\le n\) in the same window obey \[
\frac nm\le e^{1/16}<\frac{16}{15}.
\tag{FPC6}
\] Let \(\mathcal P\) be the set of all prime powers at most 264. If a
window has a nonzero coefficient, let \(m\ge3\) be its smallest
contributing prime power. Its contribution is bounded above by the full
positive sum on the exact rational superset \[
U_m=\{n\in\mathcal P:m\le n,\ 15n<16m\}.
\tag{FPC7}
\] This leaves exactly 71 cases, one for each member of
\(\mathcal P\setminus\{2\}\). A window containing no prime power has
value zero.

\subsection{Rational bounds for every individual
weight}\label{rational-bounds-for-every-individual-weight}

The companion checker uses integer arithmetic, exact fractions, and
integer square roots. It uses no floating-point value of a logarithm,
exponential, square root, or zero of zeta.

For \(0\le z\le1/3\), integrating a finite geometric expansion of
\(2/(1-z^2)\) proves \[
L_N(z):=2\sum_{j=0}^{N-1}\frac{z^{2j+1}}{2j+1}
\le \log\frac{1+z}{1-z}
\le L_N(z)+\frac{2z^{2N+1}}{(2N+1)(1-z^2)}.
\tag{FPC8}
\] Use \(N=16\). For an integer prime \(p\), put
\(k=\lfloor\log_2p\rfloor\), evaluated by its integer bit length, and \[
z_p=\frac{p-2^k}{p+2^k}\in[0,1/3).
\] Then \[
\log p=k\log2+\log\frac{1+z_p}{1-z_p}.
\tag{FPC9}
\] Apply FPC8 to both logarithms, and round the lower bound down and the
upper bound up to denominator \(10^{12}\). Denote the resulting rational
bounds by \(\lambda_p^-\) and \(\lambda_p^+\).

For each prime power \(n=p^j\), compute the integer \[
b_n=\operatorname{isqrt}(n\,10^{18}).
\] Set \(r_n^-=b_n/10^9\). If \(b_n^2=n10^{18}\), set \(r_n^+=r_n^-\);
otherwise set \(r_n^+=(b_n+1)/10^9\). By the defining
integer-square-root inequalities, \[
(r_n^-)^2\le n\le(r_n^+)^2,
\qquad
\frac{\lambda_p^-}{r_n^+}
\le\frac{\Lambda(n)}{\sqrt n}
\le\frac{\lambda_p^+}{r_n^-}.
\tag{FPC10}
\] Every denominator is positive. The checker records all 72 prime-power
indices, their prime bases, and all six rational log, square-root, and
weight bounds in the accompanying JSON.

The enumeration itself is exhaustive: trial division through the integer
square root decides every prime \(p\le264\), and repeated integer
multiplication generates every \(p^j\le264\). Every prime power has
exactly one such prime base. The complete resulting set has 72 members,
including \(243=3^5\), \(256=2^8\), 257 and 263.

\subsection{The complete finite comparison and
attainment}\label{the-complete-finite-comparison-and-attainment}

Adding the upper bounds in FPC10 on each exact set FPC7 yields the
following certified comparisons: \[
\sum_{n\in U_m}\frac{\Lambda(n)}{\sqrt n}
<\frac{1511}{1000}\quad(m\ne227),
\tag{FPC11}
\] whereas \[
U_{227}=\{227,229,233,239,241\},
\quad
\frac{1783796246}{10^9}<M_*<\frac{1783796247}{10^9}.
\tag{FPC12}
\] The lower bound in FPC12 exceeds \(1511/1000\), so this proves the
maximum over all 71 supersets exactly. These are exact rational
comparisons, not decimal estimates. Their full fractions and all
constituent indices are provided in the JSON. The full 71-case table
displays each packet and an outward-rounded upper bound with denominator
\(10^9\).

Selected packet bounds, including the packet containing 257 and 263, are

{\def\LTcaptype{none} % do not increment counter
\begin{longtable}[]{@{}rlr@{}}
\toprule\noalign{}
Minimum & Complete ratio packet & Certified upper bound \\
\midrule\noalign{}
\endhead
\bottomrule\noalign{}
\endlastfoot
227 & 227, 229, 233, 239, 241 & \(1783796247/10^9\) \\
191 & 191, 193, 197, 199 & \(1510503007/10^9\) \\
229 & 229, 233, 239, 241, 243 & \(1494205679/10^9\) \\
251 & 251, 256, 257, 263 & \(1081820789/10^9\) \\
\end{longtable}
}

The superset maximum occurs in an actual window of the original width.
Choose \[
a_*^{\rm win}=\frac{\log227+\log241}{2}.
\tag{FPC13}
\] For every prime in the displayed five-prime packet, \[
|\log p-a_*^{\rm win}|
\le\frac12\log\frac{241}{227}
<\frac12\frac{14}{227}<\frac1{32}.
\tag{FPC14}
\] Here \(\log(1+x)<x\), and \(224<227\) proves the final rational
inequality. Also \(227\cdot241<256^2\), so \(a_*^{\rm win}<\log256\).
For the lower endpoint, \(\log227>\log4=2\log2>4/3>583/800\), using the
first positive term of the logarithm series. Thus FPC13 is admissible.
Its window contains the whole five-prime packet, while FPC11--FPC12
already exclude any larger sum. It attains exactly \(M_*\), proving FPC2
and the uniform bound FPC3.

\subsection{Reproduction files}\label{reproduction-files}

\begin{itemize}
\tightlist
\item
  \texttt{check\_\allowbreak{}finite\_\allowbreak{}prime\_\allowbreak{}window\_\allowbreak{}coefficients.\allowbreak{}py}: the complete
  rational-only enumeration and all asserted comparisons.
\item
  \texttt{FINITE\_\allowbreak{}PRIME\_\allowbreak{}WINDOW\_\allowbreak{}COEFFICIENT\_\allowbreak{}CERTIFICATE.\allowbreak{}json}: every
  individual rational enclosure and every one of the 71 packet sums.
\item
  \texttt{FINITE\_\allowbreak{}PRIME\_\allowbreak{}WINDOW\_\allowbreak{}PACKETS.\allowbreak{}md}: the complete
  human-readable packet table, with rational upper bounds.
\end{itemize}

Run the Python checker in this directory. It reproduces the JSON and the
complete packet table, and fails if any stated comparison fails. The
exact maximum is the symbolic sum FPC2; its rational enclosure is FPC3.
No source window, prime-power coefficient, or endpoint has been
replaced.

\section{Exact rational packet table}\label{exact-rational-packet-table}

Each upper bound is the displayed integer divided by \(10^9\). It is an
outward rounding of the exact upper fraction in the companion JSON.

{\def\LTcaptype{none} % do not increment counter
\begin{longtable}[]{@{}
  >{\raggedleft\arraybackslash}p{(\linewidth - 4\tabcolsep) * \real{0.3333}}
  >{\raggedright\arraybackslash}p{(\linewidth - 4\tabcolsep) * \real{0.3333}}
  >{\raggedleft\arraybackslash}p{(\linewidth - 4\tabcolsep) * \real{0.3333}}@{}}
\toprule\noalign{}
\begin{minipage}[b]{\linewidth}\raggedleft
Minimum prime power
\end{minipage} & \begin{minipage}[b]{\linewidth}\raggedright
All prime powers in its ratio superset
\end{minipage} & \begin{minipage}[b]{\linewidth}\raggedleft
Upper numerator over \(10^9\)
\end{minipage} \\
\midrule\noalign{}
\endhead
\bottomrule\noalign{}
\endlastfoot
3 & 3 & 634284101 \\
4 & 4 & 346573591 \\
5 & 5 & 719762516 \\
7 & 7 & 735484905 \\
8 & 8 & 245064536 \\
9 & 9 & 366204097 \\
11 & 11 & 722992628 \\
13 & 13 & 711388957 \\
16 & 16, 17 & 860441965 \\
17 & 17 & 687155170 \\
19 & 19 & 675500630 \\
23 & 23 & 653795740 \\
25 & 25 & 321887583 \\
27 & 27 & 211428034 \\
29 & 29 & 625291138 \\
31 & 31, 32 & 739294578 \\
32 & 32 & 122532268 \\
37 & 37 & 593631249 \\
41 & 41, 43 & 1153540161 \\
43 & 43 & 573577641 \\
47 & 47, 49 & 839588912 \\
49 & 49 & 277987165 \\
53 & 53 & 545361537 \\
59 & 59, 61 & 1057193623 \\
61 & 61, 64 & 612986861 \\
64 & 64, 67 & 600328359 \\
67 & 67, 71 & 1019571991 \\
71 & 71, 73 & 1008047325 \\
73 & 73 & 502160296 \\
79 & 79, 81, 83 & 1098700093 \\
81 & 81, 83 & 607098802 \\
83 & 83 & 485030770 \\
89 & 89 & 475794504 \\
97 & 97, 101, 103 & 1380386597 \\
101 & 101, 103, 107 & 1367634487 \\
103 & 103, 107, 109 & 1357762463 \\
107 & 107, 109, 113 & 1345804284 \\
109 & 109, 113 & 894064869 \\
113 & 113 & 444715238 \\
121 & 121, 125, 127, 128 & 853061211 \\
125 & 125, 127, 128, 131 & 1061018700 \\
127 & 127, 128, 131 & 917066197 \\
128 & 128, 131 & 487214102 \\
131 & 131, 137, 139 & 1264826923 \\
137 & 137, 139 & 838878955 \\
139 & 139 & 418536617 \\
149 & 149, 151, 157 & 1221772044 \\
151 & 151, 157 & 811832789 \\
157 & 157, 163, 167 & 1198547906 \\
163 & 163, 167, 169, 173 & 1384116861 \\
167 & 167, 169, 173 & 985143554 \\
169 & 169, 173, 179 & 976824721 \\
173 & 173, 179, 181 & 1165922166 \\
179 & 179, 181 & 774124661 \\
181 & 181, 191, 193 & 1145259138 \\
191 & 191, 193, 197, 199 & 1510503007 \\
193 & 193, 197, 199 & 1130461840 \\
197 & 197, 199 & 751645111 \\
199 & 199, 211 & 743669539 \\
211 & 211, 223 & 730528112 \\
223 & 223, 227, 229, 233 & 1438337443 \\
227 & 227, 229, 233, 239, 241 & 1783796247 \\
229 & 229, 233, 239, 241, 243 & 1494205679 \\
233 & 233, 239, 241, 243 & 1135135222 \\
239 & 239, 241, 243, 251 & 1126789316 \\
241 & 241, 243, 251, 256, 257 & 1162009892 \\
243 & 243, 251, 256, 257 & 808703063 \\
251 & 251, 256, 257, 263 & 1081820789 \\
256 & 256, 257, 263 & 733057292 \\
257 & 257, 263 & 689735593 \\
263 & 263 & 343593738 \\
\end{longtable}
}

\end{document}
